% page 305 du fac-simile (image p-304 du scan)
% bloc 165.1 x 237.7 mm ; marge gauche 36.9 mm
\pgc{15.240mm}{0.000mm}{
\l{\cel{56}297}
\l{}
\l{konko. (\sou{zool.}) Envelopilo kalka di ula moluski. - EFIRS.}
\l{}
\l{konkava. Qua karakterizesas da kurveso sferala kava. - DEFIS.}
\l{}
\l{konklavo. (\sou{religio katol.}) Loko ube su inkluzas klefe la kar-}
\l{\cel{3}dinali pos la morto di papo, por exekutar la elekto di lua}
\l{\cel{3}sucedonto. - DEFIS.}
\l{}
\l{konkluzar. (\sou{trans.)}\cel{2}I. Finar, o finigar, per solvo definitiva. -}
\l{\cel{3}II. (\sou{logiko}) Deduktar la konsequi de la premisi pozita. -}
\l{\cel{3}EFIRS.}
\l{}
\l{konkoido.\cel{1}\sou{(geom}.) Kurvo ek du branchi infinita, qui havas kom}
\l{\cel{3}asimptoto rekto fixa, e simetra relate la perpendiklo trasita}
\l{\cel{3}de punto fixa a ta rekto. - DEFIS.}
\l{}
\l{\sou{konkologio}. La parto di zoologio qua koncernas la konko-moluski.}
\l{\cel{3}EFIS.}
\l{}
\l{\sou{konkordar}. (\sou{netrans., kun, pri, pro}\cel{1}ke). I. (\sou{pri plura personi}).}
\l{\cel{3}Opinionar, sentar, intersame. - II. (\sou{pri kozi}).Inter-relatar}
\l{\cel{3}harmonioze pro la existo inter li di raporto exakta. - DEFIRS.}
\l{}
\l{\sou{konkordato}. I. Interkonsento da la papo kun suvereno, pri la}
\l{\cel{3}yuri rispektiva di la Stato e di la Eklezio. - II. Kontrato}
\l{\cel{3}per qua komercisto qua ne plus povas pagar sua debaji obtenas,}
\l{\cel{3}de sua debarii, pago-fristo o la diminuto di parto de lua de-}
\l{\cel{3}bajo. - DEFIRS.}
\l{}
\l{\sou{konkreciono}. (\sou{patol.}) Agregajo solida, qua naskas acidente en}
\l{\cel{3}ula tisui. - DEFIS.}
\l{}
\l{\sou{konkreta}. Qua expresas ulo reala. - DEFIRS.}
\l{}
\l{\sou{konkubo}. I. Singla de la du personi qui kun-vivas spozatre,}
\l{\cel{3}-quankam ne intermariajita. - II. (\sou{che la Romani, antique)}.}
\l{\cel{3}Singla de la personi qui interunionis per la konkubinato.}
\l{\cel{3}- DEFIS.}
\l{}
\l{\sou{konkubinato}. (\sou{che la Romani, antique}). Sorto di mariajo, infe-}
\l{\cel{3}riora ye la mariajo civila - e la unika pri qua darfis kon-}
\l{\cel{3}tratar la ne-Romani. - DEFIS.}
\l{}
\l{\sou{konkupicenco}. (en\cel{1}\sou{la linguo kristanismala}). Inklineso, da}
\l{\cel{3}persono dekadinta, ad la plezuri sensala (precipue la sexuala)}
\l{\cel{3}- DEFIS.}
\l{}
\l{\sou{konkurencar}. (\sou{trans.}) (\sou{aludante personi qui persequas la sama}}
\l{\cel{3}\sou{skopo}).Inter-rivalesar pri sua interesti (precipue pri la}
\l{\cel{3}personi di la sama komerco-fako, industrio-fako). - DFIRS.}
\l{}
\l{\sou{konkursar}. (\sou{netrans., pri}). Esar en la rangi kun altri, pri}
\l{\cel{3}obtenor premio, nominasor, e c. - DFIRS.}
}
